
\documentclass{article}
\usepackage{colortbl}
\usepackage{makecell}
\usepackage{multirow}
\usepackage{supertabular}

\begin{document}

\newcounter{utterance}

\centering \large Interaction Transcript for game `matchit\_ascii', experiment `similar\_grid\_1', episode 1 with deepseek{-}v4{-}pro.
\vspace{24pt}

{ \footnotesize  \setcounter{utterance}{1}
\setlength{\tabcolsep}{0pt}
\begin{supertabular}{c@{$\;$}|p{.15\linewidth}@{}p{.15\linewidth}p{.15\linewidth}p{.15\linewidth}p{.15\linewidth}p{.15\linewidth}}
    \# & $\;$A & \multicolumn{4}{c}{Game Master} & $\;\:$B\\
    \hline

    \theutterance \stepcounter{utterance}  
    & & \multicolumn{4}{p{0.6\linewidth}}{
        \cellcolor[rgb]{0.9,0.9,0.9}{
            \makecell[{{p{\linewidth}}}]{
                \texttt{\tiny{[P1$\langle$GM]}}
                \texttt{Du deltar i ett samarbetsbaserat gissningsspel.} \\
\texttt{Målet är att ta reda på om detta rutnät och en annan bild som bara jag kan se är samma.} \\
\texttt{Beskriv ditt rutnät först. Sedan kommer jag att ge min beskrivning och vi kan ställa frågor till varandra om rutnäten för att ta reda på om de är samma. Börja nu din korta rutnätsbeskrivning med "BESKRIVNING:" följt av beskrivningen. Lägg inte till något annat. Det här är ditt rutnät:} \\
\texttt{X X X X X} \\
\texttt{X X X X ▢} \\
\texttt{X X X ▢ ▢} \\
\texttt{X X ▢ ▢ ▢} \\
\texttt{X ▢ ▢ ▢ ▢} \\
            }
        }
    }
    & \\ \\

    \theutterance \stepcounter{utterance}  
    & \multicolumn{4}{p{0.6\linewidth}}{
        \cellcolor[rgb]{0.9,0.9,0.9}{
            \makecell[{{p{\linewidth}}}]{
                \texttt{\tiny{[P1$\rangle$GM]}}
                \texttt{BESKRIVNING: Ett 5x5 rutnät med X och tomma rutor. Översta raden är helt fylld med X. Andra raden har fyra X och en tom ruta längst till höger. Tredje raden har tre X och två tomma rutor till höger. Fjärde raden har två X och tre tomma rutor till höger. Femte raden har ett X längst till vänster och fyra tomma rutor.} \\
            }
        }
    }
    & & \\ \\

    \theutterance \stepcounter{utterance}  
    & & & \multicolumn{2}{p{0.3\linewidth}}{
        \cellcolor[rgb]{0.9,0.9,0.9}{
            \makecell[{{p{\linewidth}}}]{
                \texttt{\tiny{[GM$|$GM]}}
                \texttt{continue} \\
            }
        }
    }
    & & \\ \\

    \theutterance \stepcounter{utterance}  
    & & \multicolumn{4}{p{0.6\linewidth}}{
        \cellcolor[rgb]{0.9,0.9,0.9}{
            \makecell[{{p{\linewidth}}}]{
                \texttt{\tiny{[GM$\rangle$P2]}}
                \texttt{Du deltar i ett samarbetsbaserat gissningsspel.} \\
\texttt{Målet är att ta reda på om detta rutnät och en annan bild som bara jag kan se är samma.} \\
\texttt{Beskriv ditt rutnät först. Sedan kommer jag att ge min beskrivning och vi kan ställa frågor till varandra om rutnäten för att ta reda på om de är samma. Börja nu din korta rutnätsbeskrivning med "BESKRIVNING:" följt av beskrivningen. Lägg inte till något annat. Det här är ditt rutnät:} \\
\texttt{X X X X X} \\
\texttt{▢ X X X X} \\
\texttt{▢ ▢ X X X} \\
\texttt{▢ ▢ ▢ X X} \\
\texttt{▢ ▢ ▢ ▢ X} \\
            }
        }
    }
    & \\ \\

    \theutterance \stepcounter{utterance}  
    & & & \multicolumn{4}{p{0.6\linewidth}}{
        \cellcolor[rgb]{0.9,0.9,0.9}{
            \makecell[{{p{\linewidth}}}]{
                \texttt{\tiny{[GM$\langle$P2]}}
                \texttt{BESKRIVNING: Ett 5x5 rutnät med X i alla celler på och ovanför diagonalen från övre vänstra till nedre högra hörnet, och tomma rutor (▢) i alla celler under diagonalen.} \\
            }
        }
    }
     \\ \\

    \theutterance \stepcounter{utterance}  
    & & & \multicolumn{2}{p{0.3\linewidth}}{
        \cellcolor[rgb]{0.9,0.9,0.9}{
            \makecell[{{p{\linewidth}}}]{
                \texttt{\tiny{[GM$|$GM]}}
                \texttt{continue} \\
            }
        }
    }
    & & \\ \\

    \theutterance \stepcounter{utterance}  
    & & \multicolumn{4}{p{0.6\linewidth}}{
        \cellcolor[rgb]{0.9,0.9,0.9}{
            \makecell[{{p{\linewidth}}}]{
                \texttt{\tiny{[GM$\rangle$P2]}}
                \texttt{Detta är min BESKRIVNINGBESKRIVNING: Ett 5x5 rutnät med X och tomma rutor. Översta raden är helt fylld med X. Andra raden har fyra X och en tom ruta längst till höger. Tredje raden har tre X och två tomma rutor till höger. Fjärde raden har två X och tre tomma rutor till höger. Femte raden har ett X längst till vänster och fyra tomma rutor.} \\
\texttt{Ställ nu en fråga för att ta reda på nya aspekter av mitt rutnät som kan skilja sig från ditt rutnät. Börja med "FRÅGA:" och lägg inte till något annat.} \\
            }
        }
    }
    & \\ \\

    \theutterance \stepcounter{utterance}  
    & & & \multicolumn{4}{p{0.6\linewidth}}{
        \cellcolor[rgb]{0.9,0.9,0.9}{
            \makecell[{{p{\linewidth}}}]{
                \texttt{\tiny{[GM$\langle$P2]}}
                \texttt{FRÅGA: I din beskrivning nämner du att andra raden har en tom ruta längst till höger. Kan du specificera om den tomma rutan är i kolumn 5, och i så fall, är de tomma rutorna på de efterföljande raderna också samlade längst till höger i stigande antal?} \\
            }
        }
    }
     \\ \\

    \theutterance \stepcounter{utterance}  
    & & & \multicolumn{2}{p{0.3\linewidth}}{
        \cellcolor[rgb]{0.9,0.9,0.9}{
            \makecell[{{p{\linewidth}}}]{
                \texttt{\tiny{[GM$|$GM]}}
                \texttt{continue} \\
            }
        }
    }
    & & \\ \\

    \theutterance \stepcounter{utterance}  
    & & \multicolumn{4}{p{0.6\linewidth}}{
        \cellcolor[rgb]{0.9,0.9,0.9}{
            \makecell[{{p{\linewidth}}}]{
                \texttt{\tiny{[P1$\langle$GM]}}
                \texttt{Detta är min BESKRIVNINGBESKRIVNING: Ett 5x5 rutnät med X i alla celler på och ovanför diagonalen från övre vänstra till nedre högra hörnet, och tomma rutor (▢) i alla celler under diagonalen.} \\
\texttt{FRÅGA: I din beskrivning nämner du att andra raden har en tom ruta längst till höger. Kan du specificera om den tomma rutan är i kolumn 5, och i så fall, är de tomma rutorna på de efterföljande raderna också samlade längst till höger i stigande antal?Börja ditt svar med "SVAR:" och lägg inte till något annat.} \\
            }
        }
    }
    & \\ \\

    \theutterance \stepcounter{utterance}  
    & \multicolumn{4}{p{0.6\linewidth}}{
        \cellcolor[rgb]{0.9,0.9,0.9}{
            \makecell[{{p{\linewidth}}}]{
                \texttt{\tiny{[P1$\rangle$GM]}}
                \texttt{SVAR: Ja, den tomma rutan på andra raden är i kolumn 5. Tredje raden har tomma rutor i kolumn 4 och 5. Fjärde raden har tomma rutor i kolumn 3, 4 och 5. Femte raden har tomma rutor i kolumn 2, 3, 4 och 5.} \\
            }
        }
    }
    & & \\ \\

    \theutterance \stepcounter{utterance}  
    & & & \multicolumn{2}{p{0.3\linewidth}}{
        \cellcolor[rgb]{0.9,0.9,0.9}{
            \makecell[{{p{\linewidth}}}]{
                \texttt{\tiny{[GM$|$GM]}}
                \texttt{continue} \\
            }
        }
    }
    & & \\ \\

    \theutterance \stepcounter{utterance}  
    & & \multicolumn{4}{p{0.6\linewidth}}{
        \cellcolor[rgb]{0.9,0.9,0.9}{
            \makecell[{{p{\linewidth}}}]{
                \texttt{\tiny{[P1$\langle$GM]}}
                \texttt{Ställ nu en fråga för att ta reda på nya aspekter av mitt rutnät som kan skilja sig från ditt rutnät. Börja med "FRÅGA:" och lägg inte till något annat.} \\
            }
        }
    }
    & \\ \\

    \theutterance \stepcounter{utterance}  
    & \multicolumn{4}{p{0.6\linewidth}}{
        \cellcolor[rgb]{0.9,0.9,0.9}{
            \makecell[{{p{\linewidth}}}]{
                \texttt{\tiny{[P1$\rangle$GM]}}
                \texttt{FRÅGA: I ditt rutnät, är den översta raden helt fylld med X i alla fem kolumner, och är den nedersta raden helt tom förutom ett X längst till vänster i kolumn 1?} \\
            }
        }
    }
    & & \\ \\

    \theutterance \stepcounter{utterance}  
    & & & \multicolumn{2}{p{0.3\linewidth}}{
        \cellcolor[rgb]{0.9,0.9,0.9}{
            \makecell[{{p{\linewidth}}}]{
                \texttt{\tiny{[GM$|$GM]}}
                \texttt{continue} \\
            }
        }
    }
    & & \\ \\

    \theutterance \stepcounter{utterance}  
    & & \multicolumn{4}{p{0.6\linewidth}}{
        \cellcolor[rgb]{0.9,0.9,0.9}{
            \makecell[{{p{\linewidth}}}]{
                \texttt{\tiny{[GM$\rangle$P2]}}
                \texttt{SVAR: Ja, den tomma rutan på andra raden är i kolumn 5. Tredje raden har tomma rutor i kolumn 4 och 5. Fjärde raden har tomma rutor i kolumn 3, 4 och 5. Femte raden har tomma rutor i kolumn 2, 3, 4 och 5.} \\
\texttt{FRÅGA: I ditt rutnät, är den översta raden helt fylld med X i alla fem kolumner, och är den nedersta raden helt tom förutom ett X längst till vänster i kolumn 1?Börja ditt svar med "SVAR:" och lägg inte till något annat.} \\
            }
        }
    }
    & \\ \\

    \theutterance \stepcounter{utterance}  
    & & & \multicolumn{4}{p{0.6\linewidth}}{
        \cellcolor[rgb]{0.9,0.9,0.9}{
            \makecell[{{p{\linewidth}}}]{
                \texttt{\tiny{[GM$\langle$P2]}}
                \texttt{SVAR: Ja, den översta raden är helt fylld med X i alla fem kolumner (kolumn 1 till 5). Den nedersta raden har ett X längst till vänster i kolumn 1 och är tom i kolumn 2, 3, 4 och 5.} \\
\\ 
\texttt{FRÅGA: Betyder det att ditt rutnät har X i precis de positioner där kolumnnumret är mindre än eller lika med radnumret (om vi räknar rader uppifrån och kolumner från vänster), och tomma rutor där kolumnnumret är större än radnumret?} \\
            }
        }
    }
     \\ \\

    \theutterance \stepcounter{utterance}  
    & & & \multicolumn{2}{p{0.3\linewidth}}{
        \cellcolor[rgb]{0.9,0.9,0.9}{
            \makecell[{{p{\linewidth}}}]{
                \texttt{\tiny{[GM$|$GM]}}
                \texttt{continue} \\
            }
        }
    }
    & & \\ \\

    \theutterance \stepcounter{utterance}  
    & & \multicolumn{4}{p{0.6\linewidth}}{
        \cellcolor[rgb]{0.9,0.9,0.9}{
            \makecell[{{p{\linewidth}}}]{
                \texttt{\tiny{[GM$\rangle$P2]}}
                \texttt{Ställ nu en fråga för att ta reda på nya aspekter av mitt rutnät som kan skilja sig från ditt rutnät. Börja med "FRÅGA:" och lägg inte till något annat.} \\
            }
        }
    }
    & \\ \\

    \theutterance \stepcounter{utterance}  
    & & & \multicolumn{4}{p{0.6\linewidth}}{
        \cellcolor[rgb]{0.9,0.9,0.9}{
            \makecell[{{p{\linewidth}}}]{
                \texttt{\tiny{[GM$\langle$P2]}}
                \texttt{FRÅGA: Har ditt rutnät någon rad där X och tomma rutor inte följer ett sammanhängande mönster, till exempel en rad där en tom ruta förekommer till vänster om ett X?} \\
            }
        }
    }
     \\ \\

    \theutterance \stepcounter{utterance}  
    & & & \multicolumn{2}{p{0.3\linewidth}}{
        \cellcolor[rgb]{0.9,0.9,0.9}{
            \makecell[{{p{\linewidth}}}]{
                \texttt{\tiny{[GM$|$GM]}}
                \texttt{continue} \\
            }
        }
    }
    & & \\ \\

    \theutterance \stepcounter{utterance}  
    & & \multicolumn{4}{p{0.6\linewidth}}{
        \cellcolor[rgb]{0.9,0.9,0.9}{
            \makecell[{{p{\linewidth}}}]{
                \texttt{\tiny{[P1$\langle$GM]}}
                \texttt{SVAR: Ja, den översta raden är helt fylld med X i alla fem kolumner (kolumn 1 till 5). Den nedersta raden har ett X längst till vänster i kolumn 1 och är tom i kolumn 2, 3, 4 och 5.} \\
\\ 
\texttt{FRÅGA: Betyder det att ditt rutnät har X i precis de positioner där kolumnnumret är mindre än eller lika med radnumret (om vi räknar rader uppifrån och kolumner från vänster), och tomma rutor där kolumnnumret är större än radnumret?} \\
\texttt{FRÅGA: Har ditt rutnät någon rad där X och tomma rutor inte följer ett sammanhängande mönster, till exempel en rad där en tom ruta förekommer till vänster om ett X?Börja ditt svar med "SVAR:" och lägg inte till något annat.} \\
            }
        }
    }
    & \\ \\

    \theutterance \stepcounter{utterance}  
    & \multicolumn{4}{p{0.6\linewidth}}{
        \cellcolor[rgb]{0.9,0.9,0.9}{
            \makecell[{{p{\linewidth}}}]{
                \texttt{\tiny{[P1$\rangle$GM]}}
                \texttt{SVAR: Ja, mitt rutnät har X i positioner där kolumnnumret är mindre än eller lika med radnumret, och tomma rutor där kolumnnumret är större än radnumret. Nej, ingen rad har en tom ruta till vänster om ett X; alla rader följer mönstret med sammanhängande X från vänster följt av sammanhängande tomma rutor till höger.} \\
            }
        }
    }
    & & \\ \\

    \theutterance \stepcounter{utterance}  
    & & & \multicolumn{2}{p{0.3\linewidth}}{
        \cellcolor[rgb]{0.9,0.9,0.9}{
            \makecell[{{p{\linewidth}}}]{
                \texttt{\tiny{[GM$|$GM]}}
                \texttt{continue} \\
            }
        }
    }
    & & \\ \\

    \theutterance \stepcounter{utterance}  
    & & \multicolumn{4}{p{0.6\linewidth}}{
        \cellcolor[rgb]{0.9,0.9,0.9}{
            \makecell[{{p{\linewidth}}}]{
                \texttt{\tiny{[P1$\langle$GM]}}
                \texttt{Ställ nu en fråga för att ta reda på nya aspekter av mitt rutnät som kan skilja sig från ditt rutnät. Börja med "FRÅGA:" och lägg inte till något annat.} \\
            }
        }
    }
    & \\ \\

    \theutterance \stepcounter{utterance}  
    & \multicolumn{4}{p{0.6\linewidth}}{
        \cellcolor[rgb]{0.9,0.9,0.9}{
            \makecell[{{p{\linewidth}}}]{
                \texttt{\tiny{[P1$\rangle$GM]}}
                \texttt{FRÅGA: I ditt rutnät, beskriver du diagonalen som går från övre vänstra till nedre högra hörnet. Är alla celler på denna diagonal fyllda med X, och är diagonalen den skiljelinje så att alla celler ovanför diagonalen också är X medan alla celler under är tomma?} \\
            }
        }
    }
    & & \\ \\

    \theutterance \stepcounter{utterance}  
    & & & \multicolumn{2}{p{0.3\linewidth}}{
        \cellcolor[rgb]{0.9,0.9,0.9}{
            \makecell[{{p{\linewidth}}}]{
                \texttt{\tiny{[GM$|$GM]}}
                \texttt{continue} \\
            }
        }
    }
    & & \\ \\

    \theutterance \stepcounter{utterance}  
    & & \multicolumn{4}{p{0.6\linewidth}}{
        \cellcolor[rgb]{0.9,0.9,0.9}{
            \makecell[{{p{\linewidth}}}]{
                \texttt{\tiny{[GM$\rangle$P2]}}
                \texttt{SVAR: Ja, mitt rutnät har X i positioner där kolumnnumret är mindre än eller lika med radnumret, och tomma rutor där kolumnnumret är större än radnumret. Nej, ingen rad har en tom ruta till vänster om ett X; alla rader följer mönstret med sammanhängande X från vänster följt av sammanhängande tomma rutor till höger.} \\
\texttt{FRÅGA: I ditt rutnät, beskriver du diagonalen som går från övre vänstra till nedre högra hörnet. Är alla celler på denna diagonal fyllda med X, och är diagonalen den skiljelinje så att alla celler ovanför diagonalen också är X medan alla celler under är tomma?Börja ditt svar med "SVAR:" och lägg inte till något annat.} \\
            }
        }
    }
    & \\ \\

    \theutterance \stepcounter{utterance}  
    & & & \multicolumn{4}{p{0.6\linewidth}}{
        \cellcolor[rgb]{0.9,0.9,0.9}{
            \makecell[{{p{\linewidth}}}]{
                \texttt{\tiny{[GM$\langle$P2]}}
                \texttt{SVAR: Ja, alla celler på diagonalen från övre vänstra till nedre högra hörnet är fyllda med X. Alla celler ovanför diagonalen är också X, och alla celler under diagonalen är tomma. Så diagonalen fungerar precis som skiljelinjen du beskriver.} \\
            }
        }
    }
     \\ \\

    \theutterance \stepcounter{utterance}  
    & & & \multicolumn{2}{p{0.3\linewidth}}{
        \cellcolor[rgb]{0.9,0.9,0.9}{
            \makecell[{{p{\linewidth}}}]{
                \texttt{\tiny{[GM$|$GM]}}
                \texttt{continue} \\
            }
        }
    }
    & & \\ \\

    \theutterance \stepcounter{utterance}  
    & & \multicolumn{4}{p{0.6\linewidth}}{
        \cellcolor[rgb]{0.9,0.9,0.9}{
            \makecell[{{p{\linewidth}}}]{
                \texttt{\tiny{[GM$\rangle$P2]}}
                \texttt{Ställ nu en fråga för att ta reda på nya aspekter av mitt rutnät som kan skilja sig från ditt rutnät. Börja med "FRÅGA:" och lägg inte till något annat.} \\
            }
        }
    }
    & \\ \\

    \theutterance \stepcounter{utterance}  
    & & & \multicolumn{4}{p{0.6\linewidth}}{
        \cellcolor[rgb]{0.9,0.9,0.9}{
            \makecell[{{p{\linewidth}}}]{
                \texttt{\tiny{[GM$\langle$P2]}}
                \texttt{FRÅGA: Finns det någon skillnad överhuvudtaget mellan våra rutnät, till exempel i antalet X totalt, eller är de identiska i både mönster och positioner för alla 25 celler?} \\
            }
        }
    }
     \\ \\

    \theutterance \stepcounter{utterance}  
    & & & \multicolumn{2}{p{0.3\linewidth}}{
        \cellcolor[rgb]{0.9,0.9,0.9}{
            \makecell[{{p{\linewidth}}}]{
                \texttt{\tiny{[GM$|$GM]}}
                \texttt{continue} \\
            }
        }
    }
    & & \\ \\

    \theutterance \stepcounter{utterance}  
    & & \multicolumn{4}{p{0.6\linewidth}}{
        \cellcolor[rgb]{0.9,0.9,0.9}{
            \makecell[{{p{\linewidth}}}]{
                \texttt{\tiny{[P1$\langle$GM]}}
                \texttt{SVAR: Ja, alla celler på diagonalen från övre vänstra till nedre högra hörnet är fyllda med X. Alla celler ovanför diagonalen är också X, och alla celler under diagonalen är tomma. Så diagonalen fungerar precis som skiljelinjen du beskriver.} \\
\texttt{FRÅGA: Finns det någon skillnad överhuvudtaget mellan våra rutnät, till exempel i antalet X totalt, eller är de identiska i både mönster och positioner för alla 25 celler?Börja ditt svar med "SVAR:" och lägg inte till något annat.} \\
            }
        }
    }
    & \\ \\

    \theutterance \stepcounter{utterance}  
    & \multicolumn{4}{p{0.6\linewidth}}{
        \cellcolor[rgb]{0.9,0.9,0.9}{
            \makecell[{{p{\linewidth}}}]{
                \texttt{\tiny{[P1$\rangle$GM]}}
                \texttt{SVAR: Baserat på dina svar verkar våra rutnät vara identiska. Antalet X är 15 i båda (5+4+3+2+1), och mönstret med X på och ovanför diagonalen och tomma rutor under är exakt samma. Positionerna för alla 25 celler överensstämmer. Jag kan inte identifiera någon skillnad mellan våra rutnät.} \\
            }
        }
    }
    & & \\ \\

    \theutterance \stepcounter{utterance}  
    & & & \multicolumn{2}{p{0.3\linewidth}}{
        \cellcolor[rgb]{0.9,0.9,0.9}{
            \makecell[{{p{\linewidth}}}]{
                \texttt{\tiny{[GM$|$GM]}}
                \texttt{continue} \\
            }
        }
    }
    & & \\ \\

    \theutterance \stepcounter{utterance}  
    & & \multicolumn{4}{p{0.6\linewidth}}{
        \cellcolor[rgb]{0.9,0.9,0.9}{
            \makecell[{{p{\linewidth}}}]{
                \texttt{\tiny{[P1$\langle$GM]}}
                \texttt{Ställ nu en fråga för att ta reda på nya aspekter av mitt rutnät som kan skilja sig från ditt rutnät. Börja med "FRÅGA:" och lägg inte till något annat.} \\
            }
        }
    }
    & \\ \\

    \theutterance \stepcounter{utterance}  
    & \multicolumn{4}{p{0.6\linewidth}}{
        \cellcolor[rgb]{0.9,0.9,0.9}{
            \makecell[{{p{\linewidth}}}]{
                \texttt{\tiny{[P1$\rangle$GM]}}
                \texttt{FRÅGA: Finns det någon cell i ditt rutnät som innehåller ett annat tecken än X eller tom ruta, eller någon avvikelse som en ram, färg, skugga eller extra markering runt specifika celler?} \\
            }
        }
    }
    & & \\ \\

    \theutterance \stepcounter{utterance}  
    & & & \multicolumn{2}{p{0.3\linewidth}}{
        \cellcolor[rgb]{0.9,0.9,0.9}{
            \makecell[{{p{\linewidth}}}]{
                \texttt{\tiny{[GM$|$GM]}}
                \texttt{continue} \\
            }
        }
    }
    & & \\ \\

    \theutterance \stepcounter{utterance}  
    & & \multicolumn{4}{p{0.6\linewidth}}{
        \cellcolor[rgb]{0.9,0.9,0.9}{
            \makecell[{{p{\linewidth}}}]{
                \texttt{\tiny{[GM$\rangle$P2]}}
                \texttt{SVAR: Baserat på dina svar verkar våra rutnät vara identiska. Antalet X är 15 i båda (5+4+3+2+1), och mönstret med X på och ovanför diagonalen och tomma rutor under är exakt samma. Positionerna för alla 25 celler överensstämmer. Jag kan inte identifiera någon skillnad mellan våra rutnät.} \\
\texttt{FRÅGA: Finns det någon cell i ditt rutnät som innehåller ett annat tecken än X eller tom ruta, eller någon avvikelse som en ram, färg, skugga eller extra markering runt specifika celler?Börja ditt svar med "SVAR:" och lägg inte till något annat.} \\
            }
        }
    }
    & \\ \\

    \theutterance \stepcounter{utterance}  
    & & & \multicolumn{4}{p{0.6\linewidth}}{
        \cellcolor[rgb]{0.9,0.9,0.9}{
            \makecell[{{p{\linewidth}}}]{
                \texttt{\tiny{[GM$\langle$P2]}}
                \texttt{SVAR: Nej, det finns inga andra tecken än X och tomma rutor i mitt rutnät. Det finns inga ramar, färger, skuggor eller extra markeringar runt några specifika celler. Alla celler är enhetliga i stil. Jag drar slutsatsen att våra rutnät är identiska.} \\
            }
        }
    }
     \\ \\

    \theutterance \stepcounter{utterance}  
    & & & \multicolumn{2}{p{0.3\linewidth}}{
        \cellcolor[rgb]{0.9,0.9,0.9}{
            \makecell[{{p{\linewidth}}}]{
                \texttt{\tiny{[GM$|$GM]}}
                \texttt{continue} \\
            }
        }
    }
    & & \\ \\

    \theutterance \stepcounter{utterance}  
    & & \multicolumn{4}{p{0.6\linewidth}}{
        \cellcolor[rgb]{0.9,0.9,0.9}{
            \makecell[{{p{\linewidth}}}]{
                \texttt{\tiny{[GM$\rangle$P2]}}
                \texttt{Fatta nu ett beslut. Vad tror du: är ditt rutnät och det andra rutnätet som beskrivs samma? Skriv "BESLUT: samma rutnät" om du tror att de är samma rutnät eller "BESLUT: olika rutnät" om du tror att de är olika rutnät. Lägg inte till något annat.} \\
            }
        }
    }
    & \\ \\

    \theutterance \stepcounter{utterance}  
    & & & \multicolumn{4}{p{0.6\linewidth}}{
        \cellcolor[rgb]{0.9,0.9,0.9}{
            \makecell[{{p{\linewidth}}}]{
                \texttt{\tiny{[GM$\langle$P2]}}
                \texttt{BESLUT: samma rutnät} \\
            }
        }
    }
     \\ \\

    \theutterance \stepcounter{utterance}  
    & & & \multicolumn{2}{p{0.3\linewidth}}{
        \cellcolor[rgb]{0.9,0.9,0.9}{
            \makecell[{{p{\linewidth}}}]{
                \texttt{\tiny{[GM$|$GM]}}
                \texttt{continue} \\
            }
        }
    }
    & & \\ \\

    \theutterance \stepcounter{utterance}  
    & & & \multicolumn{2}{p{0.3\linewidth}}{
        \cellcolor[rgb]{0.9,0.9,0.9}{
            \makecell[{{p{\linewidth}}}]{
                \texttt{\tiny{[GM$|$GM]}}
                \texttt{loss} \\
            }
        }
    }
    & & \\ \\

    \theutterance \stepcounter{utterance}  
    & & \multicolumn{4}{p{0.6\linewidth}}{
        \cellcolor[rgb]{0.9,0.9,0.9}{
            \makecell[{{p{\linewidth}}}]{
                \texttt{\tiny{[P1$\langle$GM]}}
                \texttt{SVAR: Nej, det finns inga andra tecken än X och tomma rutor i mitt rutnät. Det finns inga ramar, färger, skuggor eller extra markeringar runt några specifika celler. Alla celler är enhetliga i stil. Jag drar slutsatsen att våra rutnät är identiska.} \\
\texttt{Fatta nu ett beslut. Vad tror du: är ditt rutnät och det andra rutnätet som beskrivs samma? Skriv "BESLUT: samma rutnät" om du tror att de är samma rutnät eller "BESLUT: olika rutnät" om du tror att de är olika rutnät. Lägg inte till något annat.} \\
            }
        }
    }
    & \\ \\

    \theutterance \stepcounter{utterance}  
    & \multicolumn{4}{p{0.6\linewidth}}{
        \cellcolor[rgb]{0.9,0.9,0.9}{
            \makecell[{{p{\linewidth}}}]{
                \texttt{\tiny{[P1$\rangle$GM]}}
                \texttt{BESLUT: samma rutnät} \\
            }
        }
    }
    & & \\ \\

    \theutterance \stepcounter{utterance}  
    & & & \multicolumn{2}{p{0.3\linewidth}}{
        \cellcolor[rgb]{0.9,0.9,0.9}{
            \makecell[{{p{\linewidth}}}]{
                \texttt{\tiny{[GM$|$GM]}}
                \texttt{continue} \\
            }
        }
    }
    & & \\ \\

    \theutterance \stepcounter{utterance}  
    & & & \multicolumn{2}{p{0.3\linewidth}}{
        \cellcolor[rgb]{0.9,0.9,0.9}{
            \makecell[{{p{\linewidth}}}]{
                \texttt{\tiny{[GM$|$GM]}}
                \texttt{loss} \\
            }
        }
    }
    & & \\ \\

\end{supertabular}
}

\end{document}
